\section{Sensing and Reacting}
\label{sec:sensing_and_reacting}

% Explain how the system's sensors and actuators are integrated into your solution 
% to implement the "smart" and autonomic reactions of your system.
% Provide a brief explanation of their behaviour and include the electronic schematic 
% of the connections.
% You must also describe the most relevant parts of the code where sensor and actuator 
% signals are handled and how they interact.
% If you did not use the "real" sensors and actuators, indicate how you have simulated 
% their behaviour in your system.

% TODO: Write sensing and reacting description

\subsection{Sensors Integration}
% Explain how sensors are integrated

\subsection{Actuators Integration}
% Explain how actuators are integrated

\subsection{Smart and Autonomic Reactions}
% Describe the smart and autonomic reactions

\subsection{Electronic Schematic}
% Include the electronic schematic of the connections

% Example: Include electronic schematic
% \begin{figure}[H]
%     \centering
%     \includegraphics[width=0.8\textwidth]{images/schematic.png}
%     \caption{Electronic Schematic}
%     \label{fig:schematic}
% \end{figure}

\subsection{Relevant Code Implementation}
% Describe the most relevant parts of the code

% Example: Include code snippet
% \begin{lstlisting}[language=Python, caption={Sensor Data Processing}]
% def process_sensor_data(sensor_value):
%     # Process sensor data
%     if sensor_value > THRESHOLD:
%         activate_actuator()
%     return processed_value
% \end{lstlisting}

\subsection{Simulation (if applicable)}
% If you did not use "real" sensors and actuators, indicate how you have 
% simulated their behaviour
